\documentclass[../main.tex]{subfiles}
\begin{document}
\chapter{Verification and Validation}
\label{chap:verification}

\section{State of the Art}
Verification asks whether a computational model and its implementation correctly solve the equations they claim to solve; validation asks whether those equations adequately represent the physical system for an intended use. Both require documented evidence, but they answer different questions. A general treatment of verification and validation in scientific computing is given by Oberkampf and Roy \cite{oberkampfroy2010}; reproducible computational research also depends on recording data, code, parameters, and software context \cite{sandve2013reproducible}. This focused review uses these foundational sources, checked 2026-10-07; it is a practical checklist, not a survey of domain-specific experimental standards.

\section{Verification layers}
\subsection{Symbolic and equation checks}
Before numerical integration, verify coordinate order, dimensions, force signs, and constraint row order. For the five-link model, the symbolic residual vector includes ten pin equations and three fixed-pose equations. Differentiation produces \(J=\partial\Phi/\partial q\) and the acceleration bias \(\dot J\dot q\). These expressions are evaluated at a common state and compared against the independent numerical implementation. The symbolic ideal acceleration constraint and force-balance residuals should simplify or evaluate to zero within arithmetic precision.

The notebook applies these checks in this order:
\begin{enumerate}
\item Keep $m_i,I_i,L_i,g$ symbolic while constructing $T,V,\Phi,J$, and $\dot J\dot q$; verify that the expressions contain no floating-point coefficients and that each $L_i$ remains present in the relevant geometry.
\item Check dimensions: $q,\dot q,\ddot q\in\mathbb R^{15}$, $\Phi,\lambda,\dot J\dot q\in\mathbb R^{13}$, and the KKT matrix is $28\times28$.
\item At one common projected state, substitute numerical values into the symbolic model and compare $\Phi$, $J$, and $\dot J\dot q$ elementwise with the separately written numerical geometry routines.
\item Solve the KKT system and compute $r_f=M\ddot q+J^T\lambda-Q$ and $r_a=J\ddot q+\dot J\dot q$ for the ideal system, or the matching stabilized residual when stabilization is enabled. Report translational force rows in newtons separately from rotational moment rows in newton-metres.
\end{enumerate}
For this derivation the symbolic model reports 15 coordinates, symbolic lengths $L_1,\ldots,L_5$, ten pin rows, and three fixed-pose rows. Its exact-symbol check confirms that no floating-point coefficients have entered $\Phi$, $J$, $M$, $T$, $V$, or $Q_g$; the ideal 28-by-28 KKT system contains no Baumgarte gains, which occur only in the separate stabilized constraint. These checks ensure numerical geometry has not accidentally been baked into the symbolic equations.
For the notebook's initial-state comparison, the maximum symbolic/numeric differences are $3.331\times10^{-16}$ for $\Phi$, $2.220\times10^{-16}$ for $J$, and $4.337\times10^{-19}$ for $\dot J\dot q$. These values document that implementation pair at that pose.

\subsection{Linear solve checks}
At each acceleration solve, check the KKT residual
\begin{equation}
 r_{\mathrm{KKT}}=\begin{bmatrix}M&J^T\\J&0\end{bmatrix}
 \begin{bmatrix}\ddot q\\\lambda\end{bmatrix}-
 \begin{bmatrix}Q\\-\dot J\dot q-2\zeta\omega_BJ\dot q-\omega_B^2\Phi\end{bmatrix}.
\end{equation}
For the upper KKT block, report translational force and rotational moment residuals separately in newtons and newton-metres, or scale each row by a documented characteristic force or moment before combining them. For the lower block, report pin and grounded-translation acceleration residuals in metres per second squared separately from the grounded angular-acceleration residual in radians per second squared. Likewise, $\Phi$ contains pin-position residuals in metres and a fixed-angle residual in radians, while $J\dot q$ contains translational speeds and angular rates; report pin and ground components separately or use explicit scales. Confirm that $J$ has the expected rank at the tested pose. A successful factorization alone is not evidence that equations, signs, or units are correct.

\subsection{Trajectory checks}
Because the constraint vector contains both translational and angular rows, report separate component norms:
\begin{equation}
 \begin{aligned}
 e_{p,\mathrm{pin}}&=\max_{1\le j\le10}|\Phi_j|,&
 e_{p,\mathrm{ground}}&=\max(|\Phi_{11}|,|\Phi_{12}|),&
 e_{\theta,\mathrm{ground}}&=|\Phi_{13}|,\\
 e_{v,\mathrm{pin}}&=\max_{1\le j\le10}|(J\dot q)_j|,&
 e_{v,\mathrm{ground}}&=\max(|(J\dot q)_{11}|,|(J\dot q)_{12}|),&
 e_{\omega,\mathrm{ground}}&=|(J\dot q)_{13}|.
 \end{aligned}
\end{equation}
The pin position and velocity norms have units of metres and metres per second; the ground translation and translation-rate norms have the same respective units, while the ground angle and angular-rate norms are in radians and radians per second. The trajectory check loops over every saved state, reevaluates $\Phi(q_k)$ and $J(q_k)$, and records the maxima of these components. It also confirms that the solver reports success, reaches the requested final time, and returns finite state values. Repeat with tighter tolerances or a smaller maximum step and compare the trajectory and quantities relevant to the intended use. Output sampling density alone is not an integration convergence study.

\subsection{Validation evidence}
The current five-link model is an ideal planar rigid-body mechanism. Its notebook verifies internal consistency and numerical closure but does not compare predictions against a physical prototype or measurement data. It therefore demonstrates code verification and numerical consistency, not empirical validation. A validation study would additionally need measured geometry, masses/inertias, joint behavior, loading, sensor uncertainty, and comparison metrics appropriate to the intended application \cite{oberkampfroy2010}.

\section{Acceptance checks for the examples}
\begin{center}
\small
\begin{tabular}{@{}>{\raggedright\arraybackslash}p{0.22\textwidth}>{\raggedright\arraybackslash}p{0.38\textwidth}>{\raggedright\arraybackslash}p{0.31\textwidth}@{}}
\toprule
Evidence & Check & Five-link result \\
\midrule
Initial consistency & Pin position and velocity residuals after projection & \(4.441\times10^{-16}\,\mathrm m\); \(4.718\times10^{-16}\,\mathrm{m\,s^{-1}}\) \\
Initial grounding & Translation, angle, translation-rate, angular-rate residuals & \(0\,\mathrm m\); \(0\,\mathrm{rad}\); \(1.561\times10^{-17}\,\mathrm{m\,s^{-1}}\); \(1.648\times10^{-17}\,\mathrm{s^{-1}}\) \\
Local mobility & Rank of \(J\) at the initial pose & Rank 13; two remaining DOF \\
Equation implementation & Symbolic/numeric \(\Phi\), \(J\), and bias comparison & Differences below \(4\times10^{-16}\) \\
Force solve & Initial KKT residual by component family & \(1.998\times10^{-15}\,\mathrm N\) force; \(4.663\times10^{-15}\,\mathrm{N\,m}\) moment \\
Integration & Solver reaches 10 s successfully & 2,001 saved samples \\
Constraint preservation & Maximum pin position and velocity errors & \(1.108\times10^{-9}\,\mathrm m\); \(1.247\times10^{-7}\,\mathrm{m\,s^{-1}}\) \\
Grounding & Translation, angle, translation-rate, angular-rate maxima & \(1.333\times10^{-17}\,\mathrm m\); \(3.031\times10^{-19}\,\mathrm{rad}\); \(1.092\times10^{-15}\,\mathrm{m\,s^{-1}}\); \(1.648\times10^{-17}\,\mathrm{s^{-1}}\) \\
Energy behavior & Relative mechanical-energy drift in the baseline run & \(3.158\times10^{-7}\) \\
\bottomrule
\end{tabular}
\end{center}
These values describe one parameter set, one solver configuration, and sampled output states. They do not guarantee accuracy between samples, validate the idealized model against a physical mechanism, or establish robustness near singular configurations.

\section{Algorithmic outline}
\begin{enumerate}
 \item Verify symbolic expressions and their dimensions before assigning numerical parameters.
 \item Compare symbolic and independently coded kinematics, Jacobians, and acceleration bias at multiple admissible states.
 \item Check rank and KKT residuals at representative configurations.
 \item Integrate from projected initial data and inspect solver success and state finiteness.
 \item Evaluate position, velocity, and fixed-pose errors at every saved state.
 \item Refine tolerances and step limits, vary stabilization gains, and report trajectory and residual sensitivity rather than relying on a single run.
 \item Label any comparison with analytical solutions, benchmarks, or measurements according to the evidence actually available.
\end{enumerate}

\section{Summary}
Symbolic identities, independent implementation comparisons, linear residuals, trajectory constraints, and refinement studies provide complementary verification evidence. Physical validation requires experimental or otherwise independent evidence not present in this example.

\ifSubfilesClassLoaded{\backmatter\bookbibliography}{}
\end{document}